\corrigesChapitre{16}{Lumière : propagation et distances}
\par\noindent\begin{minipage}{\linewidth}
\begin{corrige}{1}{Source ou diffuseur}
\begin{questions}
\item Non : la lampe émet sa lumière, la feuille diffuse celle qu’elle reçoit.
\item La lampe allumée.
\end{questions}
\end{corrige}
\end{minipage}
\par\noindent\begin{minipage}{\linewidth}
\begin{corrige}{2}{Le bon sens}
\begin{questions}
\item De la feuille vers l’œil.
\item Non : la vision exige la réception de lumière.
\end{questions}
\pcFigureEx{figures/cor-vision.pdf}{La lumière reçue par l’œil vient de la feuille éclairée ; l’œil n’émet pas le rayon nécessaire à cette vision.}
\end{corrige}
\end{minipage}
\par\noindent\begin{minipage}{\linewidth}
\begin{corrige}{3}{La trajectoire}
\begin{questions}
\item En ligne droite.
\item Oui : une flèche indique le sens de propagation.
\end{questions}
\end{corrige}
\end{minipage}
\par\noindent\begin{minipage}{\linewidth}
\begin{corrige}{4}{Une grande vitesse}
\begin{questions}
\item $300\,000\unite{km/s}$.
\item Non : la propagation sur une distance finie nécessite une durée non nulle.
\end{questions}
\end{corrige}
\end{minipage}
\par\noindent\begin{minipage}{\linewidth}
\begin{corrige}{5}{La zone d’ombre}
\begin{questions}
\item L’objet intercepte une partie des trajets lumineux issus de la source.
\item Non : transmission, absorption, diffusion et géométrie du milieu changent la situation.
\end{questions}
\pcFigureEx{figures/cor-ombre.pdf}{L’objet opaque intercepte une partie des rayons issus de la source ponctuelle du modèle.}
\end{corrige}
\end{minipage}
\par\noindent\begin{minipage}{\linewidth}
\begin{corrige}{6}{Un trajet en une seconde}
\begin{questions}
\item $600\,000\unite{km}$.
\item $30\,000\unite{km}$.
\end{questions}
\end{corrige}
\end{minipage}
\par\noindent\begin{minipage}{\linewidth}
\begin{corrige}{7}{Une durée solaire}
\begin{questions}
\item $500\unite s$.
\item $8\unite{min}\,20\unite s$.
\end{questions}
\par\smallskip\noindent\textbf{Pour comprendre.}\enspace $\Delta t=d/c=150\,000\,000/300\,000=500\unite s$. Comme $500=8\times60+20$, cela représente 8 min 20 s. Il s’agit d’un trajet simple Soleil-Terre.\par
\end{corrige}
\end{minipage}
\par\noindent\begin{minipage}{\linewidth}
\begin{corrige}{8}{Un aller simple lunaire}
\begin{questions}
\item $384\,000/300\,000=1{,}28\unite s$.
\item $2{,}56\unite s$ en négligeant les mouvements pendant le trajet.
\end{questions}
\pcFigureEx{figures/cor-lune-trajets.pdf}{L’aller-retour parcourt deux fois la distance Terre-Lune.}
\end{corrige}
\end{minipage}
\par\noindent\begin{minipage}{\linewidth}
\begin{corrige}{9}{Une étoile}
\begin{questions}
\item Environ 12 ans.
\item Non : on observe l’état correspondant à la lumière partie environ 12 ans auparavant.
\end{questions}
\end{corrige}
\end{minipage}
\par\noindent\begin{minipage}{\linewidth}
\begin{corrige}{10}{Une unité à convertir}
\begin{questions}
\item $1{,}892\times10^{13}\unite{km}$.
\item Deux ans, par définition de l’année-lumière.
\end{questions}
\end{corrige}
\end{minipage}
\par\noindent\begin{minipage}{\linewidth}
\begin{corrige}{11}{Une mesure réfléchie}
\begin{questions}
\item $d=c\Delta t/2=600\unite{km}$.
\item La lumière a parcouru l’aller et le retour pendant la durée mesurée.
\end{questions}
\par\smallskip\noindent\textbf{Pour comprendre.}\enspace Distance totale : $c\Delta t=300\,000\times0{,}004=1\,200\unite{km}$. Le réflecteur est à la moitié : $d=600\unite{km}$. Les unités km/s et s donnent des km.\par
\pcFigureEx{figures/cor-echo-lumineux.pdf}{La durée mesurée couvre deux trajets de longueur d : d = cΔt/2.}
\end{corrige}
\end{minipage}
\par\noindent\begin{minipage}{\linewidth}
\begin{corrige}{12}{Une ombre qui grandit}
\begin{questions}
\item Oui dans cette configuration de source ponctuelle divergente.
\item Prolonger les rayons passant par la source et les bords de l’objet jusqu’aux deux positions de l’écran.
\end{questions}
\pcFigureEx{figures/cor-ecrans.pdf}{Prolonger les rayons passant par les bords de l’objet : l’ombre grandit quand l’écran s’éloigne.}
\end{corrige}
\end{minipage}
\par\noindent\begin{minipage}{\linewidth}
\begin{corrige}{13}{Le rayon invisible}
\begin{questions}
\item Peu de lumière est alors renvoyée latéralement vers l’œil ; le rayon dessiné n’est pas une trace matérielle visible.
\item Des particules diffusent une partie de la lumière vers l’œil et peuvent rendre le trajet perceptible, selon le dispositif.
\end{questions}
\pcFigureEx{figures/cor-diffusion.pdf}{Voir le trajet de côté demande que de la lumière soit diffusée vers l’œil.}
\end{corrige}
\end{minipage}
\par\noindent\begin{minipage}{\linewidth}
\begin{corrige}{14}{Deux échelles}
\begin{questions}
\item $\Delta t=3/(3\times10^8)=10^{-8}\unite s$.
\item Ce délai est bien inférieur à la résolution et au temps de réaction d’un chronométrage manuel.
\end{questions}
\end{corrige}
\end{minipage}
\par\noindent\begin{minipage}{\linewidth}
\begin{corrige}{15}{Une conclusion médiatique}
\begin{questions}
\item Non : son évolution dépend notamment de sa nature et de son histoire, pas uniquement du délai lumineux.
\item On reçoit une information ancienne d’un temps approximativement égal à la distance en années-lumière, exprimé en années.
\end{questions}
\end{corrige}
\end{minipage}
\clearpage
\par\noindent\begin{minipage}{\linewidth}
\phantomsection\label{cor-defi-16}\pcSolutionDefi{16}{Le miroir spatial}\begin{questions}
\item $d=300\,000\times2{,}4/2=360\,000\unite{km}$.
\item 720 000 km est la longueur totale aller-retour, pas la séparation.
\item La distance augmente si le trajet et les autres délais du dispositif sont comparables. Un retard instrumental variable doit être exclu.
\end{questions}
\pcFigureEx{figures/cor-miroir-spatial.pdf}{720 000 km est la longueur totale ; la séparation est la moitié.}
\end{minipage}
\par\noindent\begin{minipage}{\linewidth}
\phantomsection\label{cor-activite-16}\pcSousTitre{Pistes pour l’investigation}
Tracer les deux rayons limites de la petite source passant par les bords de l’objet, puis les prolonger jusqu’à l’écran. La zone située entre eux derrière l’objet est dans l’ombre. En éloignant l’écran, l’ombre grandit dans cette configuration. Une source étendue donne une pénombre et limite la netteté du modèle ponctuel.\par\end{minipage}
